\documentclass{minimal}
\setlength{\parindent}{0cm}
\setlength{\parskip}{0.2cm}

\begin{document}
\textbf{Title EN}: Semi-Local Exotic Lagrangian Tori in Dimension Four\\
\textbf{Titre FR}: Tores exotique lagrangien semi-local en dimension quatre

\textbf{Abstract EN}:
      \begin{center}
We study exotic Lagrangian tori in dimension four.
In certain Stein domains $B_{dpq}$ (which naturally appear in almost toric fibrations) we find $d+1$ families of monotone Lagrangian tori which are mutually distinct, up to symplectomorphisms.
We prove that these remain distinct under embeddings of $B_{dpq}$ into geometrically bounded symplectic four-manifolds.
We show that there are infinitely many different such embeddings when $X$ is compact and (almost) toric and hence conclude that $X$ contains arbitrarily many Lagrangian tori which are distinct up to symplectomorphisms of $X$.
In dimension four arbitrarily many different Lagrangian tori were previously known
only in del Pezzo surfaces.

Neither the embedded tori, nor the ambient space $X$ needs to be monotone for our methods to work. 
      \end{center}
      \textbf{Résumé FR}:
      \begin{center}
Nous étudions tores lagrangiens exotiques en dimension quatre.
Dans certains domaines de Stein $B_{dpq}$ (qui apparaissent naturellement dans les fibrations presque toriques), nous trouvons $d+1$ familles de tores lagrangiens monotones qui sont mutuellement distincts, jusqu'aux symplectomorphismes.
Nous prouvons que ces familles restent distinctes sous les plongements de $B_{dpq}$ dans des variétés symplectiques géométriquement bornés de quatre dimensions.
Nous montrons qu'il existe une infinité de tels plongements lorsque $X$ est compact et (presque) torique et concluons donc que $X$ contient un nombre arbitraire de tores lagrangiens distincts jusqu'aux symplectomorphismes de $X$.
En dimension quatre, un nombre arbitraire de tores lagrangiens différents était précédemment connu
seulement dans les surfaces del Pezzo.

Ni les tores plongés, ni l'espace ambiant $X$ n'ont besoin d'être monotones pour que nos méthodes fonctionnent. 
      \end{center}

\textbf{Running Title EN}: Semi-Local Exotic Lagrangian Tori in Dimension Four\\
\textbf{Titre Courant FR}: Tores exotique lagrangien semi-local en dimension quatre

\textbf{Keywords}: 

\textbf{MSC}: 53D12, 53D20, 14J17

\textbf{Authors}:\\
Joé Brendel, ETH Zürich, joe@brendel.lu\\
Johannes Hauber, Université de Neuchâtel, johannes.hauber@unine.ch\\
Joel Schmitz, Université de Neuchâtel, joel.schmitz@unine.ch
\end{document}
